\section{Controlled-Truth Validation}
\label{sec:truth}

The public city releases do not expose event identifiers, so they cannot measure
relation-recovery accuracy. We therefore separate method validation from public
auditing through a controlled generator with known ordered-event membership.
Each instance contains three required cores, eight candidate buffers, public
time intervals, and a public scalar outcome. Every core has a direct member and
one or two events receive a sequential member that can connect through the
direct member without overlapping the core. Remaining buffers are temporally
plausible distractors. Membership labels generate and evaluate the instance but
are never supplied to \methodname{} or the point rules.

We run 1,000 predeclared seeds at each $C\in\{2,3,4\}$. The composition task is
conditioned on the true selected-buffer count $q$, while hiding membership, to
isolate relation ambiguity from support-count error. The point comparator
chooses, among feasible worlds at $q$, the world maximizing an outcome-blind
temporal-proximity score. Three fixed thresholds, $0.25$, $0.50$, and $0.75$,
convert the mean selected-buffer outcome into a downstream decision. A decision
is certified only when the full frontier lies on one side of the threshold.

\begin{table}[t]
\caption{Controlled truth over 3,000 relation-incomplete event instances.
Decision error is averaged over three predeclared thresholds.}
\label{tab:truth}
\small
\centering
\begin{tabular}{@{}crrrr@{}}
\toprule
$C$ & Instances & Median width & Point MAE & Decision error \\
\midrule
2 & 1,000 & 0.446 & 0.121 & 17.9\% \\
3 & 1,000 & 0.459 & 0.118 & 17.4\% \\
4 & 1,000 & 0.473 & 0.125 & 19.0\% \\
\bottomrule
\end{tabular}
\end{table}

With the full candidate universe, the true aggregate lies in the certified
frontier in all 3,000 instances. The feasible temporal point is often close in
absolute error yet makes a threshold decision error in roughly one of six
comparisons (Table~\ref{tab:truth}). Every such error occurs when the frontier
straddles the threshold:
\[
\Pr(\text{frontier flags ambiguity}\mid \text{point decision is wrong})=1.
\]
Thus the method's downstream value is not to promise a superior missing-link
guess. It identifies when the desired conclusion is invariant to all admissible
relations and abstains precisely on these observed point-rule failures.

\paragraph{Candidate recall is not world coverage.}
We next rank buffers by the same outcome-blind temporal score and retain only
four, six, or all eight. Table~\ref{tab:truncation} reports the intermediate
six-buffer screen. Mean true-member recall remains about 84\%, but the entire
true relational world is representable only 31--33\% of the time. Nevertheless,
a nonempty frontier at the true support count exists in 98.9--100\% of
instances. Feasibility of the optimization problem is therefore not evidence
that the candidate universe still contains truth.

\begin{table}[t]
\caption{Candidate truncation after retaining six of eight buffers.}
\label{tab:truncation}
\small
\centering
\begin{tabular}{@{}crrrr@{}}
\toprule
$C$ & Member recall & World repr. & Frontier exists & Scalar covered \\
\midrule
2 & 83.4\% & 31.4\% & 98.9\% & 69.2\% \\
3 & 83.8\% & 31.9\% & 100.0\% & 71.1\% \\
4 & 83.8\% & 32.6\% & 100.0\% & 71.7\% \\
\bottomrule
\end{tabular}
\end{table}

Scalar coverage can also survive by coincidence after the joint world is
removed: this occurs in 37.8--39.2\% of instances. Consequently, validating
only an aggregate, or only average candidate recall, can conceal a misspecified
relation universe. The correct workflow reports candidate-support sensitivity
as a separate identification layer. Pair matching and connected components are
also evaluated in the artifact, but their failure rates are generator-specific:
the generator intentionally includes sequential non-clique events and bounded
occupancy. We therefore use them as diagnostic misspecification baselines, not
as estimates of real-world algorithm accuracy.
